\scanpage{10}
% archive foliation: 9
\deleted{\uline{Prop} Soient \unsure{x}\ \unsure{qui}}

On dit que $\uline{x}\in C\deleted{\unsure{(j)}}$ est \uline{acceptable} contre $j'\in J$ si \deleted{$\exists$} (posons $J'=J-\{j'\}$) \deleted{et pour $J'=J-\{j'\}$} $\exists$ $J'$-stratégie $\Sigma$ telle que $\Sigma$ soit acceptable pour $x\in C$ contre $j'$.
\margin{si $x\in C_0$ signifie donc que $x\notin G(j')$.}

Cela implique que $x\notin G(j')$. Car si on avait $x\in G(j')$, \deleted{alors} et si $\Sigma'$ était une $j'$-stratégie $x$-gagnante, alors prenons une partie maximale (à priori finie ou infinie) commençant par $x$ et compatible avec $\Sigma$ et avec $\Sigma'$, cette partie serait finie (car $\Sigma$ est gagnante pour $x$ et $j$), et on aurait (si $x_N$ est l'aboutissement) à la fois $x_N\in G(\unsure{j|j'})$ et $x_N\notin G(\unsure{j|j'})$, absurde. Mieux

\uline{Théorème 2} Soient $j'\in J$, $J'=J\setminus\{j'\}$, $x\in C$.

Pour qu'on ait $x\in\overline G(j')$ (i.e. $x\notin G(j')$), il \uline{faut} et il suffit que $x$ soit acceptable contre $j'$.
